% ---
% id: reductions/rem:chain-admission
% kind: rem
% title: "The admission rule and what this module emits"
% summary: >-
%   This module states the admission rule of a chain and decides no instance of
%   it; the protocol compiler decides both its halves. Every entry this module
%   emits carries a real constant and every artefact an empty source segment,
%   so no emitted artefact can violate the rule.
% prominence: standard
% status: accepted
% uses: [def:chain, def:attset]
% mentions: [protocol-compiler/def:chaincond, protocol-compiler/req:chain-refusals]
% assumes: []
% refs: []
% provenance: newly-derived
% lean:
% ---

\begin{remark}\label{rem:chain-admission}
This module states the admission rule of Definition~\ref{def:chain} and
decides no instance of it. It declares no general entry and emits none, every
entry it emits carrying a real scalar (Definition~\ref{def:attset}), and every
artefact it emits has an empty source segment; so no artefact of this module
can violate the rule. The rule is decided where a declaration's source
segment is assembled with its constraint list, which is the protocol
compiler's chain condition (\xref{protocol-compiler/def:chaincond}). It
decides both halves of the rule: the \emph{entry} half, a general entry
beside a non-empty source segment, and the \emph{mode} half, a declaration
carrying a non-empty source segment in untrusted-source mode, each refused by
name there (\xref{protocol-compiler/req:chain-refusals}).
\end{remark}
